\subsection{Loading, updating and deleting}

At scale 1 with asynchronous commits, HStore inserts nodes \M{ratioAingestnodes} and hyperedges \M{ratioAingestedges} as fast as HyperGraphDB. Updates differ more. Replacing node values 1,000 at a time runs \M{ratioAwriteupdate} faster in HStore (\M{hAwriteupdate} against \M{gAwriteupdate} updates per second). Deleting hyperedges is \M{ratioAchurndelete} faster and removing single members \M{ratioAchurnremove} faster. Building one hyperedge over all 50,000 nodes takes HStore a single bulk load, \M{ratioAlargeingest} faster than HyperGraphDB.

Single-update transactions are the one write workload where HStore does not lead: \M{hAlatencycommit} against \M{gAlatencycommit} commits per second (\M{ratioAlatencycommit}, ranges overlapping). With asynchronous commits, neither system waits for the disk, and the cost of one commit is dominated by work in memory. \Cref{sec:latency} shows that the two medians are close and that the difference lies in the tail.

\subsection{Reads}

HStore's largest margins are on reads that go through incidence sets. It enumerates a node's incidence set \M{ratioAreadincidence} faster than HyperGraphDB, counts the hyperedges shared by two nodes \M{ratioAreadcomembership} faster, and computes two-hop neighbourhoods \M{ratioAreadtwohop} faster. With 8 threads, incidence lookups reach \M{hAreadincidenceparallel} per second in HStore against \M{gAreadincidenceparallel} (\M{ratioAreadincidenceparallel}). When every lookup is its own read transaction (\wl{latency.read}), the ratio is \M{ratioAlatencyread}. Enumerating the members of a hyperedge, the read where HyperGraphDB's atom cache helps most, is \M{ratioAreadmembers} faster in HStore at the median, with overlapping ranges: one of the five HStore runs was slower than HyperGraphDB's median.

Two read workloads run under special conditions. Directly after a reopen, before any warm-up, HStore serves incidence lookups at \M{hAreadincidencecold} per second against \M{gAreadincidencecold} (\M{ratioAreadincidencecold}). HStore read \M{coldreadAhstore} per lookup from its files in that phase, against \M{coldreadAhgdb} for JE. After the deletes, incidence lookups are \M{ratioAreadincidencechurned} faster.

\subsection{Giant hyperedges}

On the hyperedge that contains all 50,000 nodes, HStore scans members \M{ratioAlargescan} faster (\M{hAlargescan} members per second against \M{gAlargescan}) and tests membership \M{ratioAlargeprobe} faster. A membership test in HStore is one descent of the edge's membership tree.

\subsection{Latency}
\label{sec:latency}

\begin{figure*}[t]
\centering
\input{gen/fig-latency}
\caption{Latency percentiles, median of each percentile over runs (log scale). From left: single-update commits with asynchronous and with synchronous commits, single-lookup read transactions, and commits of the open-loop writer while 7 threads read.}
\label{fig:latency}
\end{figure*}

\begin{table*}[t]
\centering\footnotesize
\caption{S1 async: latency per operation, median of each percentile over runs.}
\label{tab:latency}
\input{gen/latency-s1-async}
\end{table*}

\Cref{fig:latency} and \cref{tab:latency} give the latency distributions. For single-update commits with asynchronous commits, the medians are close, with HyperGraphDB slightly ahead: \M{latAhstorelatencycommitmed} for HStore and \M{latAhgdblatencycommitmed} for HyperGraphDB. The tails favour HStore. At p99, HStore needs \M{latAhstorelatencycommitptwo} and HyperGraphDB \M{latAhgdblatencycommitptwo}. At p99.9 the figures are \M{latAhstorelatencycommitpthree} and \M{latAhgdblatencycommitpthree}. Read transactions take \M{latAhstorelatencyreadmed} at the median in HStore and \M{latAhgdblatencyreadmed} in HyperGraphDB, and \M{latAhstorelatencyreadpthree} against \M{latAhgdblatencyreadpthree} at p99.9.

\paragraph{Reads beside a writer.} In \wl{mixed}, the writer's median rate was 1,000 updates per second in both systems; its slowest run reached \M{mixminAhstore} updates per second in HStore and \M{mixminAhgdb} in HyperGraphDB. Meanwhile the 7 readers completed \M{hAmixedread} lookups per second in HStore and \M{gAmixedread} in HyperGraphDB (\M{ratioAmixedread}). Reader latency at p99.9 was \M{latAhstoremixedreadpthree} against \M{latAhgdbmixedreadpthree}. The writer's commits, measured from their scheduled start, took \M{latAhstoremixedwriteptwo} against \M{latAhgdbmixedwriteptwo} at p99.

\paragraph{Commits during space reclamation.} In the \wl{compaction} window, the median commit latency was about the same with and without a reclamation running, for both systems. At p99, HStore's commits took \M{latAhstorecompactionptwo} during compaction against \M{latAhstorecompactionbaselineptwo} without; HyperGraphDB's took \M{latAhgdbcompactionptwo} during cleaning against \M{latAhgdbcompactionbaselineptwo} without.

\subsection{Disk footprint, deletes and reclamation}

\begin{figure}[t]
\centering
\input{gen/fig-disk}
\caption{Bytes on disk after load, after deleting 20\% of the hyperedges, and after a full space reclamation, median over runs.}
\label{fig:disk}
\end{figure}

After loading at scale 1, HStore's files take \M{hAdisk} and HyperGraphDB's \M{gAdisk} (\cref{fig:disk}). After the deletes the figures are \M{hAdiskchurned} and \M{gAdiskchurned}. HStore's files grow at this point. The deletes write new pages, and the old ones stay on disk while any of the 64 retained generations can reach them, until a later checkpoint releases them and compaction removes them. After a full reclamation and a clean shutdown, HStore is back to \M{hAdiskcompacted}, and HyperGraphDB is at \M{gAdiskcompacted}.

\subsection{Reopen and recovery}

A clean close and reopen takes \M{hAreopen} in HStore and \M{gAreopen} in HyperGraphDB (\M{ratioAreopen}). After the ingest process is killed with SIGKILL, opening the store takes \M{hArecover} and \M{gArecover} respectively, and the ranges do not overlap. HyperGraphDB opens a crashed store faster at scale 1. In every run of both systems, every node acknowledged before the kill was present after recovery.

\subsection{Resources}
\label{sec:resources}

\begin{figure}[t]
\centering
\input{gen/fig-writes}
\caption{Bytes written per operation, S1 async, median over runs (log scale). Writes include each system's log, data files, checkpoints and cleaning.}
\label{fig:writes}
\end{figure}

\begin{table*}[t]
\centering\footnotesize
\caption{Resources per configuration, median over runs. HS: HStore, HG: HyperGraphDB. Space reclaimed is the drop in disk size across the compaction window; it is negative when the writer added more than the reclamation freed.}
\label{tab:memory}
\input{gen/resources}
\end{table*}

HStore writes fewer bytes per operation than HyperGraphDB in every write workload of every configuration (\cref{fig:writes} and \cref{app:full}). A single-update commit writes \M{bpoAhstorelatencycommit} in HStore and \M{bpoAhgdblatencycommit} in HyperGraphDB. A batched update writes \M{bpoAhstorewriteupdate} against \M{bpoAhgdbwriteupdate}. Over a whole run at scale 1, HStore wrote \M{totalwrittenAhstore} and HyperGraphDB \M{totalwrittenAhgdb}.

\Cref{tab:memory} summarises memory and reclamation. HStore retained \M{retainedAhstore} of heap after ingest at scale 1, against \M{retainedAhgdb} for HyperGraphDB. Per-workload allocation and GC time are in \cref{app:full}. HStore allocates less per operation than HyperGraphDB in every read workload of every configuration. At scale 1 it allocates more per inserted node and hyperedge, and more while recovering after a crash. At scale 4 it allocates less in those workloads as well.

\subsection{Larger data: scale 4}
\label{sec:scale4}

At scale 4 the dataset has \M{incidencesC} incidences. After loading, HStore's files take \M{hCdisk} and HyperGraphDB's \M{gCdisk}, so both are larger than their 64\,MiB caches, HyperGraphDB's by a much wider margin. HStore's leads hold or grow. Hyperedge insertion becomes \M{sumCingestedges} faster (\M{hCingestedges} against \M{gCingestedges} per second), batched updates \M{sumCwriteupdate} and single-update commits \M{sumClatencycommit}. Co-membership counts are \M{sumCreadcomembership} faster, incidence lookups \M{sumCreadincidence}, and cold lookups after a reopen \M{sumCreadincidencecold}.

Member enumeration is \M{sumCreadmembers} faster: \M{hCreadmembers} hyperedges per second against HyperGraphDB's \M{gCreadmembers}. Two-hop neighbourhoods, which enumerate the members of every incident hyperedge, are \M{sumCreadtwohop} faster. During \wl{read.members}, HStore read nothing from its files and allocated \M{allocChstorereadmembers} per lookup, so every catalog leaf the probes reached was already in its node cache. A cached catalog leaf holds its page bytes and only the records that have been read (\cref{sec:systems}), which keeps the catalog within the 64\,MiB budget. HyperGraphDB serves the same lookups from its atom cache, which holds each link's targets as a Java array and is not limited by the 64\,MiB budget. It allocates \M{allocChgdbreadmembers} per member lookup. A two-hop query allocates \M{allocChstorereadtwohop} in HStore and \M{allocChgdbreadtwohop} in HyperGraphDB.

After the deletes, HStore's files grow to \M{hCdiskchurned}. The deletes rewrite a large part of the incidence index, and the replaced pages remain reachable from the 64 retained generations until the next restart releases that history. The benchmark restarts the store after the deletes, and during the 10\,s \wl{compaction.baseline} window that follows, HStore's background maintenance released and compacted that space. When the measured compaction window began, the explicit compaction took \M{compmsChstore} and found almost nothing left to reclaim, and the final size was \M{hCdiskcompacted}. JE's cleaner ran for \M{compmsChgdb} in the same window and brought HyperGraphDB from \M{gCdiskchurned} to \M{gCdiskcompacted}.

A clean reopen is \M{sumCreopen} faster in HStore, and opening after a crash takes \M{hCrecover} against \M{gCrecover}. Over the whole run, HStore wrote \M{totalwrittenChstore} and HyperGraphDB \M{totalwrittenChgdb}, most of it during cleaning. HStore's single-update commits are faster at every percentile: \M{latChstorelatencycommitptwo} against \M{latChgdblatencycommitptwo} at p99, and \M{latChstorelatencycommitpthree} against \M{latChgdblatencycommitpthree} at p99.9. The median of the per-run maxima is also lower in HStore, \M{latChstorelatencycommitmax} against \M{latChgdblatencycommitmax}, but the slowest single commit across all runs was HStore's: \M{latChstorelatencycommitworst}, against \M{latChgdblatencycommitworst} for HyperGraphDB.

\subsection{Synchronous commits}
\label{sec:sync}

With synchronous commits, every acknowledged commit has been forced to stable storage. On this machine a forced write costs about 4\,ms, and it dominates single-update transactions in both systems. HStore commits \M{hBlatencycommit} of them per second and HyperGraphDB \M{gBlatencycommit} (\M{sumBlatencycommit}). The medians are \M{latBhstorelatencycommitmed} and \M{latBhgdblatencycommitmed}, one forced write each, and HStore's lead comes from the tail: \M{latBhstorelatencycommitpthree} against \M{latBhgdblatencycommitpthree} at p99.9.

Batched work spreads the forced write over 1,000 operations. HStore inserts nodes \M{sumBingestnodes} as fast as HyperGraphDB and hyperedges \M{sumBingestedges} as fast (ranges overlapping), and updates values \M{sumBwriteupdate} as fast. The open-loop writer of \wl{mixed} commits 10 updates every 10\,ms. Both systems held that schedule in every run. Measured from their scheduled start, HStore's commits took \M{latBhstoremixedwritemed} at the median and \M{latBhstoremixedwriteptwo} at p99, and HyperGraphDB's \M{latBhgdbmixedwritemed} and \M{latBhgdbmixedwriteptwo}.

Under synchronous commits, JE reports \M{bpoBhgdblatencycommit} written per single-update commit, against \M{bpoAhgdblatencycommit} with asynchronous commits. The extra bytes come with the log flush that the adapter adds (\cref{sec:durability}). We have not established what JE writes in that flush, so the bytes-written figures for synchronous single commits in \cref{app:full} should not be read as a comparison.

Reads do not commit, and HStore leads every read workload in this configuration as well at the median, with overlapping ranges for member enumeration and giant-edge scans (\cref{tab:summary}). Recovery after SIGKILL takes \M{hBrecover} in HStore and \M{gBrecover} in HyperGraphDB at the median, and the ranges do not overlap. As in every other configuration, no run of either system lost an acknowledged commit.


